%!TEX root = ../thesis.tex
% створюємо вступ
\intro

    % Вступ є однією із самих формалізованих частин дипломної роботи. На початку 
    % ви у двох-трьох абзацах повинні окреслити проблематику та актуальність 
    % вашого дослідження, після чого переходити до мети та завдання.
    \textbf{Актуальність дослідження.}
        Стрімкий розвиток методів глибокого навчання відкрив нові можливості для розв'язання складних задач генерації та обробки візуальних даних,
        однією з яких є інтерактивна колоризація монохромних зображень.
        Актуальність розробки нових підходів у цій сфері зумовлена стабільним практичним попитом на інтелектуальні інструменти розфарбовування у креативних індустріях,
        а також для задач реставрації архівних фотографій.

        Незважаючи на значний прогрес у галузі комп'ютерного зору, існуючі архітектурні рішення мають ряд суттєвих обмежень.
        Традиційні архітектури на базі згорткової нейронної мережі (англ. Convolutional Neural Network, CNN),
        умовної генеративно-змагальної мережі (англ. Conditional Generative Adversarial Network, cGAN) та трансформера (англ. Transformer) є детермінованими.
        Такі моделі зводять задачу до однозначного відображення, що значно ускладнює природну генерацію альтернативних кольорових рішень для одного й того ж вхідного зображення без застосування специфічних модифікацій.
        З іншого боку, такі генеративні архітектури, як генеративно-змагальна мережа (англ. Generative Adversarial Network, GAN),
        дифузійна модель (англ. Diffusion Model) та авторегресивний трансформер (англ. Autoregressive Transformer), здатні забезпечувати багатоваріантність генерації.
        Однак їх значна ресурсоємність, що зумовлена високою обчислювальною складністю та тривалим часом інференсу,
        робить ці підходи важкозастосовними для миттєвої інтерактивної взаємодії у реальному часі, де від системи вимагається швидкий відгук на кольорові підказки.

        У цьому контексті перспективним напрямком є дослідження новітніх архітектур, зокрема селективних моделей простору станів (англ. Selective State Space Models, SSM).
        На сьогодні рішення на основі SSM для задач інтерактивної багатоваріантної колоризації зображень загального вигляду відсутні,
        а існуючі поодинокі дослідження обмежені вузькоспеціалізованими доменами даних.
        Оскільки чисті архітектури такого класу мають обмежені можливості щодо збереження чітких локальних просторових ознак та текстур зображення,
        виникає обґрунтована необхідність проектування саме гібридних моделей.
        Це зумовлює доцільність поєднання глобального контексту SSM із локальними властивостями згорткових підходів для створення обчислювально ефективних архітектур,
        здатних забезпечити високу якість генерації та миттєвий відгук на дії користувача.

    % є певна абстрактна недосяжна річ на кшталт загальнолюдського щастя на горизонті.
    \textbf{Мета дослідження.}  % модульний програмний комплекс
        Розробка моделі колоризації на базі запропонованої гібридної архітектури U-Net + B-Mamba,
        створення програмного забезпечення інтерактивної колоризації з можливістю вибору різних моделей
        та проведення порівняльного аналізу із сучасними базовими моделями (англ. baseline models).

    % Для досягнення мети необхідно розв'язати \textbf{задачу дослідження}, яка полягає у чомусь суттєво більш конкретному.
    % Для розв'язання задачі необхідно вирішити такі завдання:
    \textbf{Задача дослідження.}
        Для досягнення поставленої мети необхідно вирішити такі завдання:
        \begin{enumerate}
            % \item провести огляд опублікованих джерел за тематикою дослідження;
            % \item (наступний пункт, пов'язаний із теоретичним дослідженням);
            % \item (і ще один, наприклад, про експериментальну перевірку результатів);
            % \item (і взагалі, краще із науковим керівником проконсультуйтесь, як ваші завдання правильно писати).
            \item проаналізувати сучасний стан досліджень у задачі колоризації зображень та виявити обмеження існуючих методів глибокого навчання, зокрема архітектур на базі CNN, GAN, Transformer та Diffusion;
            \item сформувати набори даних для навчання і тестування, а також розробити алгоритм генерації кольорових підказок для імітації взаємодії з користувачем;
            \item спроектувати та реалізувати власну модель колоризації на базі гібридної архітектури U-Net та B-Mamba, а також інтегрувати базові моделі інших архітектур для порівняння;
            \item розробити програмне забезпечення із графічним інтерфейсом для виконання інтерактивної багатоваріантної колоризації;
            \item здійснити навчання розробленої моделі та провести комплексне оцінювання її ефективності шляхом порівняльного тестування за об'єктивними метриками і проведення користувацького дослідження.
        \end{enumerate}

    % є якісь процеси або явища загального характеру
    % (наприклад, <<інформаційні процеси в системах криптографічного захисту>>).
    \emph{Об'єктом дослідження}
        є процес інтерактивної багатоваріантної колоризації зображень.

    % є конкретний математичний чи фізичний 
    % об'єкт, який розглядається у вашій роботі та який можна трактувати
    % як певну властивість об'єкта дослідження (наприклад, <<моделі та методи
    % диференціального криптоаналізу ітеративних симетричних блочних шифрів>>).
    \emph{Предметом дослідження}
        є моделі глибокого навчання для колоризації зображень,
        метрики для їх порівняльного аналізу,
        а також методи суб'єктивного оцінювання якості інтерактивної колоризації.

    % і тут коротенький перелік (наприклад, але не обмежуючись: методи лінійної та абстрактної 
    % алгебри, теорії імовірностей, математичної статистики, комбінаторного 
    % аналізу, теорії кодування, теорії складності алгоритмів, методи 
    % комп'ютерного та статистичного моделювання) 
    При розв'язанні поставлених завдань використовувались такі \emph{методи дослідження}:
        методи глибокого навчання та методи комп'ютерного зору,
        методи об'єктивного та суб'єктивного оцінювання якості результатів.

    %  отриманих результатів полягає
    % ... -- тут необхідно 
    % перелічити, що саме нового з точки зору науки несе ваша робота. До усіх 
    % тверджень, які сюди виносяться, подумки (а іноді й явним чином) потрібно 
    % ставити слово <<вперше>> -- і ці твердження повинні залишатись істинними.
    \textbf{Наукова новизна.}
        Наукова новизна отриманих результатів полягає у наступному:
        \begin{itemize}[label=$\bullet$]
            \item вперше застосовано SSM до задачі інтерактивної багатоваріантної колоризації монохромних зображень загального вигляду,
                що дозволило розширити сферу використання цього класу моделей за межі аналізу вузькоспеціалізованих даних, а саме інфрачервоних та супутникових знімків;
            \item запропоновано нову модифіковану гібридну архітектуру, на базі поєднання U-Net та B-Mamba, для задач інтерактивної колоризації.
                За рахунок інтеграції механізму двонаправленого сканування зі спільними вагами вперше вдалося ефективно об'єднати здатність SSM
                до захоплення глобального контексту з локальною просторовою точністю згорткових мереж, суттєво оптимізувавши споживання відеопам'яті.
        \end{itemize}

    % результатів полягає... -- тут необхідно 
    % зазначити практичну користь від результатів вашої роботи. Що саме можна 
    % покращити, підвищити (або знизити), зробити гарного (або уникнути 
    % поганого) після вашого дослідження.
    \textbf{Практичне значення.}
        Практичне значення отриманих результатів полягає в тому, що розроблена обчислювально ефективна гібридна архітектура та створений на її основі програмний інструмент можуть бути використані в задачах інтерактивної колоризації зображень у творчих індустріях та реставраційних проєктах. Зокрема:
        \begin{itemize}[label=$\bullet$]
            \item завдяки оптимізації використання відеопам'яті, запропонована модель дозволяє виконувати інтерактивну багатоваріантну колоризацію на локальних системах, оснащених CUDA-сумісними графічними прискорювачами, без необхідності залучення серверних потужностей;
            \item створено програмний прототип із графічним інтерфейсом, що забезпечує можливість інтерактивного керування процесом колоризації за допомогою кольорових підказок у режимі реального часу;
            \item розроблений модульний пайплайн навчання та метричного оцінювання дозволяє розробникам швидко адаптувати запропоновану архітектуру під нові набори даних або суміжні задачі комп'ютерного зору без необхідності розробки інфраструктури з нуля.
        \end{itemize}

    % \textbf{Апробація результатів та публікації.}
        % Наприкінці вступу необхідно 
        % зазначити перелік конференцій, семінарів та публікацій, в яких викладено 
        % результати вашої роботи. Якщо результати вашої роботи ніде не 
        % доповідались, опускайте даний абзац.